%\input{../../../../Config/configPoly}
\input{configPolyLocal}
\lhead{PGE1 Dijon - S6}
\newcommand{\Var}{\operatorname{Var}}

\begin{document}

\begin{titlepage}
\thispagestyle{empty}
\begin{center}
\vspace*{2cm}
{\Huge\bfseries Probabilités et statistiques\par}
\vspace{0.5cm}
{\Large PGE1 - Semestre 6\par}
\vspace{0.8cm}
{\Large Année 2025-2026\par}
\vspace{1.5cm}
\vspace{0.5cm}
\vspace{2cm}
\begin{tabular}{rl}
Enseignant : & Maxime Berger \\
Formation : & PGE1 Dijon \\
Volume : & 12 h de CM et 5 TD \\
\end{tabular}
\vfill
\begin{minipage}{0.85\textwidth}
\centering
Ce polycopié rassemble les définitions, les outils, les méthodes de calcul et les exemples utiles pour le module. 
Les TD sont indispensables pour la préparation de l'examen.
\end{minipage}
\vfill
{\large ESTP\par}
\end{center}
\end{titlepage}

\pagenumbering{roman}
\tableofcontents

\newpage
\pagenumbering{arabic}

\chapter{Introduction}
L'objectif du module n'est pas d'accumuler des formules, mais d'acquérir une méthode : partir d'une situation réelle, identifier les bonnes variables, expliciter les hypothèses, effectuer un calcul avec la loi de probabilité adéquate, puis interpréter le résultat.

\section*{Objectifs du module}

À l'issue du cours, vous devez être capable de :
\begin{itemize}
    \item traduire une situation concrète en langage probabiliste 
    \item choisir un outil adapté parmi les lois et méthodes usuelles 
    \item calculer une probabilité, une espérance, une variance ou un intervalle de confiance 
    \item commenter le résultat en français simple dans un contexte donné 
    \item discuter la pertinence d'un modèle et les limites d'une conclusion.
\end{itemize}

\section*{Esprit du poly}

Ce document est un \textbf{support d'usage}. Il contient :
\begin{itemize}
    \item les définitions à connaître 
    \item les formules utiles 
    \item des méthodes de résolution 
    \item des exemples appliqués au domaine de la construction et du contrôle
\end{itemize}

\section*{Fils rouges}

Tout au long du poly, plusieurs situations reviennent régulièrement :
\begin{itemize}
    \item le \textbf{contrôle qualité du béton} et des éprouvettes ;
    \item la \textbf{planification de chantier} sous aléa météorologique ;
    \item le \textbf{contrôle de production} sur des lots de pièces ;
    \item l'\textbf{estimation statistique} à partir d'un échantillon de mesures.
\end{itemize}


\chapter{Modéliser l'aléa en ingénierie}
\input{chapitres/1-EspaceProbabilise}

\chapter{Conditionner l'information et décider}
\input{chapitres/2-ProbabilitesConditionnelles}

\chapter{Variables aléatoires discrètes}
\input{chapitres/4-VariablesAleatoiresDiscretes}

\chapter{Variables aléatoires continues}
\input{chapitres/5-VariablesAleatoiresContinues}

\chapter{Échantillonnage et théorème central limite}
% Chapitre 5 : Échantillonnage et TCL

\section{Population, échantillon, statistique}

\begin{Def}
La \textbf{population} est l'ensemble complet sur lequel on voudrait de l'information.
\end{Def}

\begin{Def}
Un \textbf{échantillon} est un ensemble fini d'observations prélevées dans la population.
\end{Def}

\begin{Ex}
\textbf{Exemples}
\begin{itemize}
    \item Population : toutes les éprouvettes produites par une centrale pendant une campagne.
    \item Échantillon : les $n=30$ éprouvettes contrôlées cette semaine.
\end{itemize}
\end{Ex}

\begin{Def}
Une \textbf{statistique} est une quantité calculée à partir de l'échantillon : moyenne, proportion, variance empirique, maximum, etc.
\end{Def}

\section{Moyenne et proportion empiriques}

\begin{Def}
Si l'on observe $x_1,\dots,x_n$, la \textbf{moyenne empirique} est
$$
\overline{x}=\frac{x_1+\cdots+x_n}{n}.
$$
\end{Def}

\begin{Def}
Si l'on compte un nombre $k$ de succès dans un échantillon de taille $n$, la \textbf{proportion empirique} est
$$
\widehat{p}=\frac{k}{n}.
$$
\end{Def}

\begin{Rmq}
La moyenne empirique et la proportion empirique sont des observations calculées sur l'échantillon. Elles servent à approcher la moyenne ou la proportion de la population.
\end{Rmq}

\section{Dispersion dans l'échantillon}

\begin{Def}
La \textbf{variance empirique} est
$$
s^2=\frac{1}{n}\sum_{i=1}^n (x_i-\overline{x})^2.
$$
\end{Def}

\begin{Def}
L'\textbf{écart-type empirique} est $s=\sqrt{s^2}$.
\end{Def}

\begin{Rmq}
Dans ce cours, l'objectif principal de $s$ est d'évaluer la dispersion d'un échantillon et de comprendre la précision d'une moyenne.
\end{Rmq}

\section{Moyenne d'échantillon comme variable aléatoire}

\begin{Prop}
Si $X_1,\dots,X_n$ sont indépendantes et de même loi, de moyenne $\mu$ et de variance $\sigma^2$, alors la moyenne d'échantillon
$$
\overline{X}=\frac{X_1+\cdots+X_n}{n}
$$
verifie
$$
\mathbb{E}[\overline{X}]=\mu,
\qquad
\Var(\overline{X})=\frac{\sigma^2}{n}.
$$
\end{Prop}

\begin{Rmq}
Quand la taille d'échantillon augmente, la variance de la moyenne diminue. Une moyenne calculée sur beaucoup d'observations est donc plus stable.
\end{Rmq}

\section{Théorème central limite}

\begin{Prop}[Théorème central limite]
Si $X_1,\dots,X_n$ sont indépendantes et de même loi, de moyenne $\mu$ et d'écart-type $\sigma$, alors pour $n$ grand :
$$
\frac{\overline{X}-\mu}{\sigma/\sqrt{n}}
$$
est approximativement distribuée selon la loi normale centrée réduite $\mathcal{N}(0,1)$.
\end{Prop}

\begin{Rmq}
Le théorème central limite explique pourquoi la loi normale apparaît si souvent en statistique, même lorsque la population de départ n'est pas normale.
\end{Rmq}

\section{Cas d'une proportion}

\begin{Prop}
Si l'on observe $n$ essais indépendants avec probabilité de succès $p$, alors la proportion empirique $\widehat{P}$ vérifie :
$$
\mathbb{E}[\widehat{P}] = p,
\qquad
\Var(\widehat{P}) = \frac{p(1-p)}{n}.
$$
\end{Prop}

\begin{Prop}
Pour $n$ suffisamment grand, on a l'approximation normale
$$
\frac{\widehat{P}-p}{\sqrt{p(1-p)/n}} \approx \mathcal{N}(0,1).
$$
\end{Prop}

\section{Quand l'approximation normale est-elle pertinente ?}

\begin{Meth}
Dans ce cours, on utilise l'approximation normale :
\begin{itemize}
    \item pour une moyenne d'échantillon quand $n$ est suffisamment grand ;
    \item pour une proportion quand $np$ et $n(1-p)$ ne sont pas trop petits.
\end{itemize}

En pratique, dans les exercices, l'énoncé ou le contexte donne en général les conditions pour l'utiliser.
\end{Meth}

\section{Interprétation pratique}

\begin{Ex}
\textbf{Contrôle d'un lot de béton}

On mesure la résistance de $n=36$ éprouvettes. Si la moyenne vraie est $\mu$ et l'écart-type $\sigma$, alors :
$$
\overline{X} \approx \mathcal{N}\left(\mu,\frac{\sigma^2}{36}\right).
$$

L'écart-type de la moyenne vaut $\sigma/6$. La moyenne d'échantillon est donc beaucoup moins dispersée que la résistance individuelle.
\end{Ex}

\section{Erreurs fréquentes}

\begin{itemize}
    \item Confondre la dispersion d'une observation individuelle et celle d'une moyenne.
    \item Oublier le facteur $\sqrt{n}$ dans les formules.
    \item Utiliser le TCL sans justifier que l'on travaille sur un échantillon suffisamment grand.
    \item Confondre proportion empirique $\widehat{p}$ et probabilité théorique $p$.
\end{itemize}

\section{Ce qu'il faut savoir faire}

\begin{itemize}
    \item calculer une moyenne et une proportion empiriques ;
    \item interpréter un écart-type empirique ;
    \item reconnaître la loi approximative d'une moyenne d'échantillon ;
    \item reconnaître la loi approximative d'une proportion empirique ;
    \item utiliser le TCL comme outil de calcul et non comme résultat à démontrer.
\end{itemize}


\chapter{Estimation et intervalles de confiance}
% Chapitre 6 : Estimation et intervalles de confiance

\section{Pourquoi faire de l'estimation ?}

En pratique, on ne connaît pas la moyenne ou la proportion exacte d'une population. On observe un échantillon et l'on cherche à \textbf{estimer} le paramètre inconnu.

\begin{Ex}
\begin{itemize}
    \item Estimer la résistance moyenne d'une production de béton.
    \item Estimer la proportion de pièces défectueuses sur une chaîne.
    \item Estimer la part de jours défavorables à un type d'intervention.
\end{itemize}
\end{Ex}

\section{Estimateur ponctuel}

\begin{Def}
Un \textbf{estimateur ponctuel} est une formule qui fournit une valeur unique pour approcher un paramètre inconnu.
\end{Def}

\begin{Ex}
\begin{itemize}
    \item La moyenne empirique $\overline{X}$ estime la moyenne $\mu$.
    \item La proportion empirique $\widehat{P}$ estime la proportion $p$.
\end{itemize}
\end{Ex}

\begin{Rmq}
Une estimation ponctuelle est simple à utiliser, mais elle ne donne pas à elle seule la précision de l'estimation.
\end{Rmq}

\section{Intervalle de confiance : idée générale}

\begin{Def}
Un \textbf{intervalle de confiance} est un intervalle aléatoire construit à partir de l'échantillon, censé encadrer le paramètre inconnu avec un niveau de confiance donné.
\end{Def}

\begin{Rmq}
Dans ce cours, on utilisera principalement les niveaux de confiance :
\begin{itemize}
    \item $90\%$ ;
    \item $95\%$ ;
    \item $99\%$.
\end{itemize}
\end{Rmq}

\section{Intervalle de confiance d'une moyenne}

\begin{Prop}
Si la moyenne d'échantillon est approximativement normale, alors un intervalle de confiance asymptotique de niveau $95\%$ pour $\mu$ est
$$
\left[\overline{x}-1.96\frac{\sigma}{\sqrt{n}},\ \overline{x}+1.96\frac{\sigma}{\sqrt{n}}\right]
$$
si $\sigma$ est connu.
\end{Prop}

\begin{Rmq}
Dans de nombreux exercices de ce cours, on remplace $\sigma$ par une valeur fournie par l'énoncé ou par un écart-type estimé lorsque cela est explicitement demandé.
\end{Rmq}

\begin{Ex}
On observe une résistance moyenne de $\overline{x}=32.4$ MPa sur $n=36$ éprouvettes, avec un écart-type connu $\sigma=3$ MPa.

Alors :
$$
IC_{95\%} =
\left[
32.4 - 1.96\times \frac{3}{6},
32.4 + 1.96\times \frac{3}{6}
\right]
= [31.42,\ 33.38].
$$
\end{Ex}

\section{Intervalle de confiance d'une proportion}

\begin{Prop}
Pour un grand échantillon, un intervalle de confiance asymptotique de niveau $95\%$ pour une proportion $p$ est
$$
\left[
\widehat{p} - 1.96\sqrt{\frac{\widehat{p}(1-\widehat{p})}{n}},
\widehat{p} + 1.96\sqrt{\frac{\widehat{p}(1-\widehat{p})}{n}}
\right].
$$
\end{Prop}

\begin{Ex}
Sur $200$ pièces contrôlées, $18$ sont non conformes. Alors
$$
\widehat{p}=\frac{18}{200}=0.09.
$$
Un intervalle de confiance à $95\%$ s'obtient avec
$$
0.09 \pm 1.96 \sqrt{\frac{0.09\times 0.91}{200}}.
$$
\end{Ex}

\section{Marge d'erreur}

\begin{Def}
La \textbf{marge d'erreur} est la demi-largeur de l'intervalle de confiance.
\end{Def}

\begin{Rmq}
La marge d'erreur :
\begin{itemize}
    \item diminue quand $n$ augmente ;
    \item augmente quand le niveau de confiance augmente ;
    \item dépend aussi de la dispersion des données.
\end{itemize}
\end{Rmq}

\section{Comment interpréter un intervalle ?}

\begin{Meth}
Pour commenter un intervalle de confiance :
\begin{enumerate}
    \item on rappelle le paramètre estimé ;
    \item on donne l'intervalle numérique ;
    \item on le compare à une valeur de référence ou à un seuil technique ;
    \item on conclut sans sur-interpréter.
\end{enumerate}
\end{Meth}

\begin{Ex}
Si un seuil de conformité moyen est fixé à $30$ MPa et que l'intervalle à $95\%$ pour la moyenne est $[31.42 ; 33.38]$, alors le seuil $30$ MPa est en dessous de tout l'intervalle. Les données sont donc compatibles avec une moyenne supérieure à $30$ MPa.
\end{Ex}

\section{Ce qu'un intervalle ne dit pas}

\begin{Rmq}
Un intervalle de confiance ne dit pas :
\begin{itemize}
    \item que le paramètre a $95\%$ de chance d'être dans cet intervalle une fois l'intervalle calculé ;
    \item que $95\%$ des observations sont dans l'intervalle ;
    \item qu'il n'existe aucun risque d'erreur de décision.
\end{itemize}
\end{Rmq}

\section{Erreurs fréquentes}

\begin{itemize}
    \item Oublier le facteur $\sqrt{n}$.
    \item Utiliser $p$ à la place de $\widehat{p}$ dans un exercice d'estimation de proportion.
    \item Confondre intervalle sur le paramètre et intervalle contenant les observations.
    \item Donner un résultat numérique sans interprétation.
\end{itemize}

\section{Ce qu'il faut savoir faire}

\begin{itemize}
    \item calculer une estimation ponctuelle de moyenne ou de proportion ;
    \item construire un intervalle de confiance usuel ;
    \item identifier la marge d'erreur ;
    \item conclure par rapport à un seuil de décision.
\end{itemize}


\chapter{Estimateurs et tests d'hypothèses}
% Chapitre 8 : Estimateurs et tests d'hypothèses

\section{Pourquoi parler d'estimateurs et de tests ?}

En statistique appliquée, l'ingénieur ne cherche pas seulement à calculer un intervalle de confiance. Il doit aussi :
\begin{itemize}
    \item choisir un \textbf{estimateur} pertinent pour approcher un paramètre inconnu ;
    \item décider si des données sont \textbf{compatibles} avec une hypothèse de départ ;
    \item formuler une conclusion claire pour une décision technique.
\end{itemize}

\begin{Rmq}
Ce chapitre doit être lu comme un chapitre de \textbf{prise de décision raisonnée}. On ne demande pas aux étudiants de mémoriser une théorie sophistiquée des tests, mais de comprendre ce que signifie : poser une hypothèse, mesurer un écart, puis conclure sans sur-interpréter.
\end{Rmq}

\begin{Rmq}
Dans le cadre du module `2025-2026`, ce chapitre est proposé comme un \textbf{complément de culture statistique}. Il peut nourrir la compréhension générale du cours, mais les tests d'hypothèses ne feront pas l'objet de questions spécifiques à l'examen.
\end{Rmq}

\begin{Ex}
Quelques questions typiques :
\begin{itemize}
    \item la résistance moyenne d'un béton est-elle conforme à la valeur cible ?
    \item la proportion de pièces défectueuses dépasse-t-elle le seuil acceptable ?
    \item un nouveau procédé semble-t-il meilleur que l'ancien ?
\end{itemize}
\end{Ex}

\section{Qu'est-ce qu'un estimateur ?}

\begin{Def}
Un \textbf{estimateur} est une variable aléatoire construite à partir d'un échantillon et destinée à approcher un paramètre inconnu.
\end{Def}

\begin{Ex}
\begin{itemize}
    \item la moyenne empirique $\overline{X}$ estime la moyenne $\mu$ ;
    \item la proportion empirique $\widehat{p}$ estime la proportion $p$ ;
    \item la variance empirique estime la dispersion de la population.
\end{itemize}
\end{Ex}

\begin{Rmq}
Il faut distinguer :
\begin{itemize}
    \item l'\textbf{estimateur}, qui est une formule aléatoire ;
    \item l'\textbf{estimation}, qui est la valeur numérique obtenue sur l'échantillon observé.
\end{itemize}
\end{Rmq}

\section{Qualités attendues d'un estimateur}

\begin{Def}
Un estimateur est dit \textbf{sans biais} si sa moyenne est égale au paramètre qu'il estime.
\end{Def}

\begin{Ex}
La moyenne empirique $\overline{X}$ est un estimateur sans biais de $\mu$.
\end{Ex}

\begin{Def}
Un estimateur est \textbf{précis} lorsque sa dispersion est faible.
\end{Def}

\begin{Rmq}
Dans ce cours, on retiendra surtout l'idée suivante :
\begin{itemize}
    \item un bon estimateur vise juste ;
    \item un bon estimateur varie peu d'un échantillon à l'autre.
\end{itemize}
\end{Rmq}

\section{Principe d'un test d'hypothèses}

\begin{Def}
Un \textbf{test d'hypothèses} est une procédure qui permet de confronter une hypothèse de départ, appelée \textbf{hypothèse nulle} et notée $H_0$, à des données observées.
\end{Def}

\begin{Def}
L'hypothèse alternative, notée $H_1$, décrit ce que l'on conclut si les données sont trop peu compatibles avec $H_0$.
\end{Def}

\begin{Ex}
Dans un contrôle de conformité sur une moyenne :
\begin{itemize}
    \item $H_0$ : la moyenne vaut la valeur cible ;
    \item $H_1$ : la moyenne est inférieure à la valeur cible.
\end{itemize}
\end{Ex}

\section{Les étapes d'un test}

\begin{Meth}
Pour mener un test d'hypothèses, on suit toujours le même schéma :
\begin{enumerate}
    \item écrire clairement $H_0$ et $H_1$ ;
    \item choisir une statistique de test ;
    \item calculer sa valeur sur l'échantillon ;
    \item comparer cette valeur à un seuil, ou calculer une \textit{p-value} ;
    \item conclure dans le contexte.
\end{enumerate}
\end{Meth}

\begin{Meth}
Formulation attendue dans une copie :
\begin{itemize}
    \item \textbf{étape 1} : écrire clairement ce que l'on cherche à vérifier ;
    \item \textbf{étape 2} : traduire cette question en $H_0$ et $H_1$ ;
    \item \textbf{étape 3} : faire le calcul ;
    \item \textbf{étape 4} : conclure en français simple pour l'ingénieur.
\end{itemize}
\end{Meth}

\section{Risque de première espèce}

\begin{Def}
Le \textbf{risque de première espèce}, noté $\alpha$, est la probabilité de rejeter $H_0$ alors qu'elle est vraie.
\end{Def}

\begin{Rmq}
Dans les exercices usuels, on prend souvent :
\begin{itemize}
    \item $\alpha = 5\%$ ;
    \item parfois $\alpha = 1\%$.
\end{itemize}
\end{Rmq}

\section{Lien avec la p-value}

\begin{Def}
La \textbf{p-value} est le plus petit niveau de risque pour lequel les données conduiraient à rejeter $H_0$.
\end{Def}

\begin{Meth}
Règle de décision avec la p-value :
\begin{itemize}
    \item si $\text{p-value} \leq \alpha$, on rejette $H_0$ ;
    \item sinon, on ne rejette pas $H_0$.
\end{itemize}
\end{Meth}

\begin{Rmq}
Ne pas rejeter $H_0$ ne signifie pas que $H_0$ est vraie avec certitude. Cela signifie seulement que l'échantillon ne fournit pas assez d'arguments pour la rejeter.
\end{Rmq}

\section{Test sur une moyenne : idée générale}

\begin{Ex}
\textbf{Résistance moyenne d'un béton}

On veut vérifier si une production atteint une résistance moyenne cible de $30$ MPa.

On peut poser :
\begin{itemize}
    \item $H_0 : \mu = 30$ ;
    \item $H_1 : \mu < 30$.
\end{itemize}

Si l'on observe une moyenne d'échantillon très inférieure à $30$, on sera conduit à rejeter $H_0$.
\end{Ex}

\begin{Prop}
Lorsque l'on connaît l'écart-type $\sigma$ et que la taille d'échantillon est suffisante, on utilise souvent la statistique
$$
Z=\frac{\overline{X}-\mu_0}{\sigma/\sqrt{n}}.
$$
\end{Prop}

\section{Test sur une proportion : idée générale}

\begin{Ex}
\textbf{Taux de pièces défectueuses}

On veut vérifier si la proportion de pièces défectueuses dépasse $5\%$.

On peut poser :
\begin{itemize}
    \item $H_0 : p = 0.05$ ;
    \item $H_1 : p > 0.05$.
\end{itemize}

Si la proportion observée $\widehat{p}$ est nettement supérieure à $0.05$, on s'oriente vers le rejet de $H_0$.
\end{Ex}

\begin{Prop}
Pour une proportion, la statistique usuelle est
$$
Z=\frac{\widehat{p}-p_0}{\sqrt{p_0(1-p_0)/n}}
$$
lorsque l'approximation normale est justifiée.
\end{Prop}

\section{Lien entre test et intervalle de confiance}

\begin{Rmq}
Les tests et les intervalles de confiance racontent souvent la même chose sous deux formes différentes :
\begin{itemize}
    \item si une valeur de référence est \textbf{hors} de l'intervalle de confiance, on a tendance à rejeter l'hypothèse correspondante ;
    \item si elle est \textbf{dans} l'intervalle, on ne la rejette généralement pas.
\end{itemize}
\end{Rmq}

\section{Erreurs fréquentes}

\begin{itemize}
    \item confondre \og ne pas rejeter $H_0$ \fg{} et \og prouver $H_0$ \fg ;
    \item oublier de formuler explicitement $H_0$ et $H_1$ ;
    \item conclure sans revenir au contexte technique ;
    \item oublier que le niveau $\alpha$ est un risque choisi à l'avance ;
    \item transformer un test en recette automatique sans discuter ce que mesurent réellement les données.
\end{itemize}

\section{Ce qu'il faut savoir faire}

\begin{itemize}
    \item distinguer estimateur et estimation ;
    \item citer les qualités attendues d'un estimateur ;
    \item écrire correctement $H_0$ et $H_1$ ;
    \item expliquer le rôle du risque $\alpha$ et de la p-value ;
    \item interpréter un test comme un outil d'aide à la décision ;
    \item distinguer un écart \textbf{statistiquement significatif} d'un écart \textbf{techniquement important}.
\end{itemize}


\appendix
\chapter{Fiches méthode}
% Annexe : fiches méthode

\section{Choisir un modèle probabiliste}

\begin{Meth}
Avant tout calcul, poser les quatre questions suivantes :
\begin{enumerate}
    \item Quelle est la grandeur ou le phénomène observé ?
    \item Est-ce un comptage ou une mesure ?
    \item Quelles hypothèses de simplification accepte-t-on ?
    \item Quelle est la question finale : calcul, comparaison, décision ?
\end{enumerate}
\end{Meth}

\section{Fiche réflexe : événements et conditionnement}

\begin{itemize}
    \item Traduire l'énoncé en événements nommés clairement.
    \item Si l'énoncé parle de \og sachant \fg{}, penser à la probabilité conditionnelle.
    \item Si l'énoncé comporte plusieurs cas exclusifs, penser à une partition puis aux probabilités totales.
    \item Si l'on inverse une information, penser a Bayes.
\end{itemize}

\section{Fiche réflexe : lois discrètes}

\begin{center}
\renewcommand{\arraystretch}{1.5}
\begin{tabular}{|m{3cm}|m{4.5cm}|m{4.8cm}|}
\hline
\textbf{Loi} & \textbf{Quand l'utiliser ?} & \textbf{Paramètres et grandeurs utiles} \\
\hline
Bernoulli & Un seul essai, succès ou échec & $p$, moyenne $p$, variance $p(1-p)$ \\
\hline
Binomiale & $n$ essais indépendants, même probabilité de succès & $n$, $p$, moyenne $np$, variance $np(1-p)$ \\
\hline
Géométrique & Rang du premier succès dans des essais indépendants & $p$, moyenne $\dfrac{1}{p}$ \\
\hline
Poisson & Comptage d'événements rares sur une durée ou une zone & $\lambda$, moyenne $\lambda$, variance $\lambda$ \\
\hline
\end{tabular}
\end{center}

\section{Fiche réflexe : lois continues}

\begin{center}
\renewcommand{\arraystretch}{1.5}
\begin{tabular}{|m{3cm}|m{4.5cm}|m{4.8cm}|}
\hline
\textbf{Loi} & \textbf{Quand l'utiliser ?} & \textbf{Paramètres et grandeurs utiles} \\
\hline
Uniforme & Toutes les valeurs d'un intervalle sont équiprobables & $[a,b]$, moyenne $\dfrac{a+b}{2}$ \\
\hline
Exponentielle & Temps d'attente ou durée de vie simplifiée & $\lambda$, moyenne $\dfrac{1}{\lambda}$ \\
\hline
Normale & Mesures influencées par de nombreux petits effets & $\mu$, $\sigma$, centrage-réduction \\
\hline
\end{tabular}
\end{center}

\section{Fiche réflexe : loi normale}

Si $X\sim \mathcal{N}(\mu,\sigma^2)$, on pose
$$
Z=\frac{X-\mu}{\sigma}.
$$

Alors :
\begin{itemize}
    \item $\mathbb{P}(X\leq a)=\mathbb{P}\left(Z\leq \dfrac{a-\mu}{\sigma}\right)$ ;
    \item $\mathbb{P}(a\leq X\leq b)=\mathbb{P}\left(\dfrac{a-\mu}{\sigma}\leq Z \leq \dfrac{b-\mu}{\sigma}\right)$.
\end{itemize}

Valeurs usuelles à connaître :
\begin{itemize}
    \item $1.645$ pour $90\%$ ;
    \item $1.96$ pour $95\%$ ;
    \item $2.576$ pour $99\%$.
\end{itemize}

\section{Fiche réflexe : intervalles de confiance}

\begin{itemize}
    \item IC d'une moyenne :
    $$
    \overline{x} \pm z_{\alpha/2}\frac{\sigma}{\sqrt{n}}.
    $$
    \item IC d'une proportion :
    $$
    \widehat{p} \pm z_{\alpha/2}\sqrt{\frac{\widehat{p}(1-\widehat{p})}{n}}.
    $$
\end{itemize}

\section{Checklist de rédaction}

\begin{itemize}
    \item Nommer la variable ou les événements.
    \item Justifier le choix du modèle.
    \item Écrire la formule avant l'application numérique.
    \item Donner une unité quand il y en a une.
    \item Conclure par une phrase interprétable par un ingénieur.
\end{itemize}

\section{Erreurs classiques à relire avant l'examen}

\begin{itemize}
    \item Confondre $\mathbb{P}(A|B)$ et $\mathbb{P}(B|A)$.
    \item Confondre variable discrète et variable continue.
    \item Oublier le facteur $\sqrt{n}$ dans les calculs statistiques.
    \item Confondre moyenne théorique et moyenne observée.
    \item Écrire une conclusion sans revenir à la question du contexte.
\end{itemize}


\end{document}
